\chapter{Background and related work}\label{ch:background}

\section{Concepts a reader needs}
\subsection{How PostgreSQL reports time}
\texttt{EXPLAIN (ANALYZE, FORMAT JSON)} executes the query and returns a tree of plan nodes. Each node
carries \texttt{Actual Total Time} (the time to produce all its rows, \emph{per loop}),
\texttt{Actual Loops}, \texttt{Actual Rows}, and buffer counters; the document root carries
\texttt{Execution Time} \citep{pg:explain}. Three details drive the design of CETP.
\begin{enumerate}
\item \textbf{Times are inclusive.} A node's time contains its children's time. The node's own work (its
  \emph{self-time}) is its inclusive time minus its children's inclusive time.
\item \textbf{Times are per loop.} A node executed $L$ times (for example the inner side of a nested loop)
  reports the \emph{average} time of one execution, so total inclusive time is \texttt{Actual Total Time}
  $\times L$.
\item \textbf{Parallel nodes report averages across participants.} Below a \texttt{Gather} node, the
  leader and workers each run the subtree; the reported per-node time is averaged, so node self-times sum to
  roughly the total CPU time, not wall time. I measured the ratio on the three real fixtures at
  $2.97\times$ (q19) and $3.00\times$ (q1) \evRUN\ (\cref{sec:worked}); the source comment says $2.85\times$.
\end{enumerate}

\subsection{Why parallelism dominates cross-machine differences}
PostgreSQL launches up to \texttt{max\_parallel\_workers\_per\_gather} workers (default 2, giving a
degree of parallelism of 3 with the leader). If a container is limited by a CPU quota (the Linux CFS
scheduler, \texttt{docker --cpus}), three processes time-share fewer cores and the \emph{parallel
speed-up} falls below 1: measured 0.769 at one CPU and 2.64--2.70 at two or more (\cref{tab:speedup}).
A scalar capability ratio cannot represent this; it affects only the nodes below a \texttt{Gather}.

\subsection{Metrics}\label{sec:metrics}
Let $y_i$ be the measured production runtime and $\hat y_i$ the prediction. Because
$\hat y_i = T^{\text{dev}}_i\,\hat s_i$ and $y_i=T^{\text{dev}}_i\,s_i$ share the same development time,
relative errors on runtime and on the scaling factor $s$ are identical, and the evaluation scripts score
$s$ directly.
\[
\MAPE=\frac{100}{n}\sum_i\frac{|\hat y_i-y_i|}{y_i},\qquad
\text{MdAPE}=100\cdot\operatorname{median}_i\frac{|\hat y_i-y_i|}{y_i},\qquad
q_i=\max\!\Big(\frac{\hat y_i}{y_i},\frac{y_i}{\hat y_i}\Big).
\]
MAPE is dominated by over-predictions (no upper limit) and is asymmetric; the repository therefore also
reports MdAPE and $q$-error \citep{moerkotte:qerror}. All differences quoted as significant use a
\emph{cluster} bootstrap: whole machines (folds) are resampled with replacement 2{,}000 times
(seed 42), so correlated pairs from the same machine are not counted as independent evidence
\citep{efron:bootstrap}.

\subsection{Cross-validation protocols}
\begin{center}\small
\begin{tabularx}{\textwidth}{@{}L{2.4cm}Y@{}}
\toprule \thead{Protocol} & \thead{What is held out}\\ \midrule
LOMO & one machine: tests generalisation to unseen hardware (the claim that matters)\\
LOQO & one query: tests generalisation to unseen workloads\\
LOMQO & one machine \emph{and} one query: both unseen at once\\
\bottomrule
\end{tabularx}
\end{center}
A leakage-free implementation refits \emph{every} learned component (classifier, quantile models) inside
each training fold; the repository audits this explicitly (\src{results/canonical\_leakage\_audit.csv}).
Leakage is a known source of inflated estimates in applied machine learning \citep{kaufman:leakage}.

\section{Related work}
CETP sits between three literatures.

\paragraph{Learned query performance prediction.}
Early work predicts execution time from plan or query features with regression or kernel methods
\citep{ganapathi:multiple,akdere:learning}; later work asks whether the optimizer's own cost model is
usable after calibration \citep{wu:costmodels}. Plan-structured neural networks \citep{marcus:planstructured}
and tree-based end-to-end estimators \citep{sun:endtoend} model the plan as a tree. These approaches
train and test on the \emph{same} system and hardware. CETP's setting differs: the target
\emph{machine} is the unseen variable. The most related line is zero-shot cost models
\citep{hilprecht:zeroshot}, which generalise across \emph{databases} using transferable plan features. CETP
instead transfers across \emph{hardware}, with an explicit physical model rather than a learned
representation, and with a tiny training set (five or eight machines).

\paragraph{Performance modelling and hardware characterisation.}
The roofline model \citep{williams:roofline} bounds throughput by $\min(\text{peak compute},\,\text{bandwidth}\times\text{intensity})$.
The STREAM triad \citep{mccalpin:stream} measures sustainable memory bandwidth. CETP's Gen-1
signature (bandwidth plus multi-core GFLOPS) borrows these two measurements but replaces the minimum
with a fixed-weight blend; the repository's own finding is that the weighting does not matter. The
Gen-2 probe instead times PostgreSQL's own operators, which sidesteps the argument that a
floating-point signature does not describe a database workload.

\paragraph{Cardinality and optimizer evaluation.}
\citet{leis:howgood} showed that cardinality errors, not the cost model, dominate plan quality. CETP
avoids cardinality estimation by taking measured node times from \texttt{EXPLAIN ANALYZE} and, in its
implemented-but-unvalidated cardinality term, only scales by a row ratio.

\paragraph{Uncertainty and selective prediction.}
Split conformal prediction gives distribution-free marginal coverage under exchangeability
\citep{angelopoulos:gentle}; quantile regression with gradient boosting \citep{friedman:gbm} gives
conditional estimates without that guarantee. The risk--coverage view of abstention follows
\citet{geifman:selective}. CETP uses all three and reports where each fails (\cref{sec:conformal}).

\paragraph{Measurement methodology.}
Cloud runtime measurements vary substantially between instances and over time
\citep{schad:cloudvariance,leitner:chaos}, and small environmental differences can change benchmark
conclusions \citep{mytkowicz:wrongdata}. The warm-up discard, repetition counts and CV thresholds used in
\cref{ch:impl} follow this concern.

\paragraph{What is, as far as I can tell, new here.}
I did not find prior work that (i) isolates \texttt{Gather}-subtree parallelism as the dominant
cross-machine term for PostgreSQL scaling, (ii) reports that a roofline-style signature fails across
instruction sets, and (iii) reports the associated negative results with cluster-bootstrap intervals. This
is a statement about my limited literature search (one session, six papers verified by search, the rest
from memory), not a novelty proof; a thesis-level claim needs a systematic review.
